\chapter{Análisis y discusión}

Los ensayos demostraron que una trayectoria puede producir avance repetible y
mantener inclinaciones acotadas sin ejecutar el patrón de apoyo supuesto. Esta
diferencia es relevante porque el desplazamiento visual, por sí solo, habría
permitido aceptar incorrectamente el gateo. La publicación sincronizada de
fase y contacto permitió contrastar la intención del controlador con la
interacción pie--suelo.

El debounce no mejoró mecánicamente la marcha. Su función fue exigir
persistencia antes de declarar vuelo. Por ello redujo la coincidencia temporal
global y añadió aproximadamente el retardo configurado a los despegues
delanteros. En las patas traseras, los eventos crudos duraron alrededor de
0,074~s y desaparecieron al aplicar el criterio de 0,12~s. La interpretación
defendible es que estos eventos correspondieron a interrupciones breves del
contacto de Gazebo, no a una fase de vuelo sostenida.

La corrección del temporizador sí produjo una mejora verificable de ejecución:
la cadencia pasó de ciclos cercanos a 4,80~s a ciclos de 4,32~s coherentes con
la configuración. Debido a que esta modificación cambió la velocidad real, las
bolsas anteriores y posteriores se analizaron como versiones experimentales
separadas.

La comparación Gazebo--MuJoCo utilizó finalmente el mismo contrato de pose,
IMU, contacto y estado articular. Esto permitió contrastar todas las variables
sin imputarlas, pero reveló diferencias materiales: MuJoCo avanzó mucho menos y
no produjo los despegues previstos. Se demostró equivalencia de interfaz y
procedimiento, no equivalencia dinámica ni física.

Los márgenes de estabilidad y las posiciones de pie dependen todavía de una
geometría y centro de masa nominales. Asimismo, los modelos de contacto no
representan directamente los futuros sensores físicos. Los resultados de
simulación respaldan el desarrollo del controlador y el protocolo, pero no
sustituyen la validación del robot real.

La capa correctiva se entrenó con cinco semillas y se conectó a Gazebo, pero
ninguna escala residual positiva cumplió simultáneamente los criterios de
avance y postura. La escala cero reprodujo el nominal, lo que descartó un fallo
del bypass, mientras las escalas positivas mostraron interferencia con la
marcha. En consecuencia, el resultado válido es negativo: la política se
rechaza para transferencia. La comparación fue descriptiva y no completamente
emparejada; no sustituye la campaña final prevista después de rediseñar el
entrenamiento.

Las pruebas provocadas demostraron que el supervisor recibe estímulos de pose,
contacto, estabilidad y trayectoria y ordena \texttt{stand}. Sin embargo,
\texttt{stand} es una orden lógica, no un corte de energía. En hardware deberá
existir una ruta independiente por OE, además de un rearme explícito y una
medición de falsos positivos antes de habilitar por defecto los watchdogs de
contacto, margen y telemetría.

\section{De la simulación al robot físico}

El cambio que hizo posible ejecutar PPO en el robot no fue el algoritmo, que se
mantuvo, sino el entorno de entrenamiento. Las primeras políticas se rechazaron
porque su escala correctiva interfería con la marcha; las posteriores se
entrenaron con la marcha que ejecuta el robot, una observación limitada a lo que
el robot mide y un modelo del actuador identificado con sus propios sensores.
Con esa formulación, las correcciones aprendidas fueron pequeñas y compatibles
con los servos, y la política de gateo redujo el error de seguimiento y la
desviación de las coxas frente a la marcha nominal de forma repetida.

La comparación entre robot y simulación mostró que los mayores errores no
estaban en la dinámica, sino en la interfaz entre ambos: el sentido invertido de
dos tibias, una postura base con las rodillas $\approx48^\circ$ más rectas que en
el modelo, topes que impedían alcanzar la postura pedida y una escala de los
sensores que cambió con la postura. Ninguno de estos errores podía compensarse
entrenando más, porque la política habría aprendido sobre un robot distinto del
real. Las cámaras resultaron decisivas para detectarlos, porque ofrecieron una
medida independiente de los sensores articulares.

\section{Fallos de instrumentación y validez de los resultados}

Varios ensayos que inicialmente parecían un empeoramiento del robot tenían otra
causa. Dos trotes con seguimiento bajo coincidieron con cientos de pasos de
control retrasados por la carga del PC; repetidos con el PC libre, el
seguimiento fue normal. Saltos simultáneos en todos los potenciómetros se
debieron a una masa común degradada, y varias interrupciones, a fallos del bus
I$^2$C. Si estos casos se hubieran tratado como resultados de las políticas, las
conclusiones habrían sido erróneas. Por ello cada caminata pasó una comprobación
de salud, los ensayos inválidos se conservaron con su causa documentada, y la
temporización se incorporó como una variable más del diagnóstico. A la vez, la
degradación observada al final de una jornada de unas doce horas, sin carga del
PC, indica que la fatiga de los servos es un factor real que conviene limitar
en campañas largas.

\section{Alcance de la evaluación en el robot suspendido}

La recompensa de las políticas premia avanzar recto sin inclinar el cuerpo,
pero con el robot suspendido no hay avance ni inclinación que medir. La métrica
física disponible, el seguimiento de las órdenes y la estabilidad de las coxas,
es necesaria para una buena locomoción pero no suficiente: una política podría
seguir bien sus órdenes y aun así caminar peor sobre el suelo. Las mejoras
obtenidas deben leerse, por tanto, como evidencia de que las correcciones
aprendidas son ejecutables por el hardware y no lo perturban, y no como una
demostración de mejor locomoción. La diferencia entre días en el nivel de error
(del orden de $0{,}5^\circ$ en las coxas) refuerza que las comparaciones deben
hacerse dentro de una misma sesión, con orden alternado y varias repeticiones.

\section{Elección de marchas}\label{sec:eleccion-marchas}

El paso, con las coxas fijas, resultó en la práctica un gateo continuo: una pata
en el aire cada vez, sin el desplazamiento lateral del cuerpo que lo distinguía
en el diseño. Por eso se añadió el trote, de pares diagonales, que sí es
una marcha distinta, y más tarde la carrera y el galope. El paso se retomó después con el
actuador medido, pero no debe contarse como una marcha realmente distinta del gateo. El trote resultó más exigente para el hardware: sus picos
de corriente coincidieron con el único reinicio de la interfaz de E/S y con la
mayor parte de los problemas de seguimiento, lo que hace de la alimentación de
los servos un requisito previo para ensayarlo sobre el suelo.

\section{Qué significan la convergencia y la viabilidad}

La convergencia de la recompensa no equivale a una política mejor para el robot. En
las cinco marchas, mientras la recompensa se estabilizaba, el avance por episodio
bajaba (sección~\ref{sec:convergencia-marchas}): la recompensa total sube sobre todo
porque bajan los castigos. Esto explica un resultado llamativo de la
Tabla~\ref{tab:viabilidad}: tres de las políticas elegidas tienen solo 0,2~M de pasos,
mucho antes de converger. Las copias tempranas conservan el avance de la marcha nominal;
las tardías derivan menos, pero avanzan menos y dejan de cumplir el criterio. Por
ejemplo, en el paso v24 semilla 4109 la copia de 0,2~M avanza 1,16~m con 0,07~m de
deriva, y la final, 0,87~m con 0,03~m. El gateo necesitó lo contrario, un peso de avance de 40 y 4,1~M de pasos,
para no perder avance. El peso de avance decide, por tanto, cuánto de cada cosa se
obtiene, y no hay un valor único válido para todas las marchas.

También hay que matizar lo que PPO aporta en simulación. Ninguna política se cayó,
pero la marcha nominal tampoco. La mejora de las políticas elegidas es de pocos puntos
de avance o de menos deriva, según la marcha, no un cambio cualitativo. La viabilidad
indica que la política no empeora a la marcha nominal cuando cambian las condiciones;
no que sea mucho mejor que ella.

\section{Parecido entre el robot y la simulación}

El parecido del 66 al 83~\% (Tabla~\ref{tab:parecido}) se debe a la amplitud, no a la
forma: las correlaciones casi siempre superan 0,86, de modo que el modelo del actuador
reproduce bien cuándo se mueve cada articulación, pero no cuánto. Además, la diferencia
crece en los ciclos cortos, ya que la carrera y el galope, de ciclo de 0,8~s, se mueven hasta 1,8 veces lo simulado, mientras que el gateo, el más lento, llega al 90~\% sin la tibia de sensor averiado.
Esto importa, porque una política entrenada con una amplitud distinta de la real
corrige sobre un movimiento diferente. El modelo del actuador se identificó sobre todo
con marchas lentas, por lo que conviene repetir la identificación con la carrera y el
galope antes de ensayarlas apoyadas.

Corregir una sola escala de sensor cambió el parecido en 7 puntos, lo que confirma lo
visto el 21 de octubre: buena parte de la brecha entre simulación y robot está en
la interfaz y no en la dinámica. Por otra parte, en estos ensayos la simulación
también está colgada, de modo que el parecido no dice nada sobre el contacto con el
suelo, donde se espera la mayor diferencia.
